\chapter{Gérer les chercheurs, les traders et les ingénieurs}\label{ch:fm:managing-researchers-traders-and-engineers}

Les signaux de trois chercheurs ont alimenté une même prévision combinée qui a rapporté 30~millions de \$ en un an. Interrogé sur ce que chacun y avait apporté, le desk mesurait
ce que la prévision aurait perdu sans chaque signal à tour de rôle~: 2.5~millions de \$, 2.7~millions de \$ et 4.6~millions de \$, soit moins d'un tiers du profit
cumulé des trois. Deux de ces signaux étaient deux voies vers le même effet, si bien que chacun paraissait presque sans valeur dès que l'autre était là. Le reste des
30~millions de \$ devaient être répartis autrement, et la façon dont cela a été fait a décidé lequel des trois est resté.

\section{Choisir des projets~: un portefeuille de recherche}

Un desk quantitatif consacre une grande partie de son budget à des travaux dont la valeur est incertaine~: nouveaux signaux, nouveaux marchés, nouveaux modèles, meilleure exécution. Le
responsable du desk choisit entre les projets comme un investisseur choisit entre des actifs, avec les mêmes deux questions~: quel rendement chacun doit-il attendre, et comment
les rendements se combinent.

\begin{definition}[Portefeuille de recherche, tableau de bord de projet]\label{def:fm:managing-researchers-traders-and-engineers:portfolio}
Le \emph{portefeuille de recherche}\index{portefeuille de recherche} d'un desk est l'ensemble des projets de recherche et de développement qu'il finance sur une période, choisis sous un
budget de moyens financiers et humains. Un \emph{tableau de bord de projet}\index{tableau de bord de projet} recense pour chaque candidat son coût, sa probabilité de succès, sa valeur en cas
de réussite (la valeur actuelle du P\&L qu'il apporte, nette de ses coûts de fonctionnement), son recouvrement avec ce que le desk fait déjà tourner, et les éléments
qui étayent chaque estimation.
\end{definition}

La valeur nette espérée d'un projet est $pV(1-\text{recouvrement})-c$. Le desk du chapitre compte douze candidats (à titre illustratif, en millions de \$)~: un nouveau
signal sur contrats à terme (coût 1.5, probabilité 0.30, valeur 12), des variables de flux d'options (2.0, 0.20, 15), un modèle d'exécution (1.0, 0.60, 4), un essai de données alternatives
(2.5, 0.10, 20), une deuxième variante de momentum qui recouvre à 70~\% le portefeuille en place (1.0, 0.40, 6), et sept autres. Avec un budget de 10~millions de \$, en choisissant
gloutonnement selon la valeur nette espérée par dollar de coût, il finance six projets dont la valeur nette espérée atteint 8.15~millions de \$. Sur 20\,000 années simulées,
le portefeuille retenu rapporte en moyenne 8.08~millions de \$ et perd de l'argent dans 28.5~\% des années~; un portefeuille tiré au hasard dans le même budget
rapporte 3.65~millions de \$ et perd dans 44.8~\% des années (voir la \cref{fig:fm:managing-researchers-traders-and-engineers:portfolio}).

\begin{omfigure}
\begin{tikzpicture}
\begin{axis}[width=10.5cm,height=5.2cm,ybar,bar width=6pt,xlabel={\small valeur nette réalisée des projets de l'année (millions de \$)},
  ylabel={\small part des années (\%)},xmin=-13,xmax=41,ymin=0,ymax=20,tick label style={font=\footnotesize},grid=major,grid style={gray!20},
  table/col sep=comma,legend cell align=left,
  legend style={font=\scriptsize,at={(0.5,-0.3)},anchor=north,/tikz/every even column/.append style={column sep=8pt}},legend columns=2]
\addplot[fill=omDef!55,draw=omDef] table[x=centre,y=greedy]{figdata/desk/09-managing-researchers-traders-and-engineers/portfolio.csv};
\addplot[fill=gray!40,draw=gray] table[x=centre,y=random]{figdata/desk/09-managing-researchers-traders-and-engineers/portfolio.csv};
\legend{choisi selon la valeur espérée par dollar,choisi au hasard}
\end{axis}
\end{tikzpicture}

\omcaption{Distribution de la valeur nette réalisée en un an par douze projets candidats, avec un budget de 10~millions de \$, choisis de façon gloutonne selon la valeur nette espérée par dollar
et au hasard~; 20\,000 années simulées, intervalles de 4~millions de \$. La plupart des années de recherche perdent de l'argent sur certains projets~; un bon portefeuille
perd moins souvent. Données~: \texttt{fm\_research.portfolio\_samples}.}\label{fig:fm:managing-researchers-traders-and-engineers:portfolio}
\end{omfigure}

\begin{remark}[Les entrées fragiles du tableau de bord]\label{rem:fm:managing-researchers-traders-and-engineers:inputs}
Les probabilités de succès sont les chiffres les plus fragiles d'un tableau de bord, et les plus importants. Elles viennent de l'historique propre du desk (combien de
ses projets passés de chaque type ont atteint la production, ce que consigne le journal de recherche du livre~7, chapitre~1), et elles sont optimistes quand
elles viennent du proposant. Tenir cet historique, c'est ce qui fait passer le portefeuille du débat à l'estimation.
\end{remark}

\section{Des revues sans théâtre}

Une revue de recherche (livre~7, chapitre~1) est le moment où un projet confronte ses éléments probants. Les revues tournent au théâtre lorsque leur issue ne dépend pas
de ce qui est présenté~: lorsque le parrain du projet occupe un poste élevé, lorsque le backtest est montré sans le nombre de variantes essayées, et quand un critère d'abandon
n'a jamais été rédigé. Les procédures qui préviennent ce cas figurent déjà dans la série~: pré-enregistrement de l'hypothèse et du test, un nombre d'essais
tenu dans le journal de recherche, jalons de validation assortis de critères d'abandon convenus avant l'arrivée des éléments.

\begin{method}[Une revue qui tranche]\label{met:fm:managing-researchers-traders-and-engineers:review}
\begin{enumerate}
\item Avant les travaux~: consigner l'hypothèse, le test, le critère d'abandon et le jalon de validation dans le journal de recherche.
\item À la revue~: présenter le résultat au regard du test pré-enregistré, avec le nombre d'essais réalisés et le ratio de Sharpe déflaté
(livre~4, chapitre~12).
\item Décider au jalon de validation~: poursuivre, arrêter ou changer l'hypothèse, ce qui ouvre un nouveau nombre d'essais.
\item Évaluer les évaluateurs une fois par an~: comparer leurs décisions à ce qui a atteint la production et à ce qu'il a rapporté.
\end{enumerate}
\end{method}

Le nombre d'essais compte plus que tout autre chiffre de la revue. Un backtest affichant un ratio de Sharpe de 1.2 sur cinq ans de données quotidiennes présente un
ratio de Sharpe probabiliste de 0.996~: pris seul, il est presque certainement réel. S'il est le meilleur de 40 variantes testées, le meilleur attendu de 40
variantes sans valeur (avec une dispersion des ratios de Sharpe de 0.5 par an) est de 1.09, et le ratio de Sharpe déflaté tombe à 0.59
(voir la \cref{fig:fm:managing-researchers-traders-and-engineers:dsr}). Le même résultat, après 200 variantes, relève plus vraisemblablement de la chance que du talent.

\begin{omfigure}
\begin{tikzpicture}
\begin{axis}[width=10.5cm,height=5cm,xmode=log,log basis x=10,xlabel={\small nombre de variantes testées (échelle logarithmique)},
  ylabel={\small ratio de Sharpe déflaté},xmin=1,xmax=500,ymin=0,ymax=1.05,tick label style={font=\footnotesize},grid=major,grid style={gray!20},
  table/col sep=comma,xtick={1,10,100},xticklabels={1,10,100}]
\addplot[omDef,thick,mark=*,mark size=1.3pt] table[x=trials,y=dsr]{figdata/desk/09-managing-researchers-traders-and-engineers/dsr.csv};
\draw[omThm,dashed] (axis cs:1,0.5) -- (axis cs:500,0.5);
\end{axis}
\end{tikzpicture}

\omcaption{Probabilité qu'un backtest quotidien sur cinq ans, avec un ratio de Sharpe annuel de 1.2, reflète un ratio de Sharpe réel supérieur au meilleur
attendu par le seul hasard parmi les variantes testées, en fonction de leur nombre (dispersion des ratios de Sharpe de 0.5 par an). Données~:
\texttt{fm\_research.dsr\_curve}, sur le \texttt{firm.multitest} du livre~4.}\label{fig:fm:managing-researchers-traders-and-engineers:dsr}
\end{omfigure}

\section{Qui a fait gagner de l'argent~: attribuer le mérite}

\begin{definition}[Contribution par exclusion (leave-one-out), attribution du mérite]\label{def:fm:managing-researchers-traders-and-engineers:credit}
La \emph{contribution par exclusion}\index{contribution par exclusion} (leave-one-out) d'un contributeur à un résultat combiné est le résultat obtenu avec tous les contributeurs
moins le résultat sans ce contributeur. L'\emph{attribution du mérite}\index{attribution du mérite} est une règle qui répartit un résultat combiné entre
les contributeurs à ce résultat, et sur laquelle reposent leur reconnaissance et leur rémunération.
\end{definition}

Une prévision combinée est un jeu au sens de la valeur de Shapley (livre~12, chapitre~6)~: chaque ensemble $S$ de contributeurs a une valeur $v(S)$,
le P\&L de la meilleure combinaison de leurs signaux. La \omterm{def:fm:managing-researchers-traders-and-engineers:credit}{contribution par exclusion}, $v(N)-v(N\setminus\{i\})$, mesure ce que chacun apporte en dernier~; lorsque deux
contributeurs apportent la même information, chacun n'apporte que peu en dernier et tous deux paraissent sans valeur. Une attribution dépendant de l'ordre donne à chaque
contributeur ce qu'il a apporté à son arrivée, si bien que le premier de deux signaux qui se recouvrent reçoit tout et le second rien. La valeur de Shapley fait la moyenne
sur tous les ordres possibles.

\begin{proposition}[Ce que garantit le partage de Shapley]\label{prop:fm:managing-researchers-traders-and-engineers:shapley}
La valeur de Shapley
$\phi_i=\frac{1}{n!}\sum_{\pi}\bigl(v(P_i^\pi\cup\{i\})-v(P_i^\pi)\bigr)$, la moyenne sur les ordres $\pi$ de ce que $i$ apporte à l'ensemble $P_i^\pi$ des
contributeurs qui le précèdent, vérifie $\sum_i\phi_i=v(N)$~; attribue un mérite égal aux contributeurs qui apportent la même chose à tout ensemble~; et n'attribue rien à un
contributeur qui n'apporte rien à aucun ensemble. La \omterm{def:fm:managing-researchers-traders-and-engineers:credit}{contribution par exclusion} vérifie les deux dernières mais, en général, pas la première~: lorsque les contributions sont
substituables, $v(N)-v(N\setminus\{i\})$ sommé sur $i$ est inférieur à $v(N)$.
\end{proposition}

\begin{proof}[Preuve partielle]
Pour chaque ordre, les valeurs apportées se télescopent en $v(N)-v(\emptyset)=v(N)$~; la moyenne conserve la somme. Les contributeurs symétriques ont
des distributions identiques de valeurs apportées selon l'ordre~; un contributeur qui n'apporte rien n'apporte rien dans aucun ordre. Pour la \omterm{def:fm:managing-researchers-traders-and-engineers:credit}{contribution par exclusion}, prenons deux
contributeurs identiques~: chacun n'apporte rien en dernier, si bien que tous deux reçoivent zéro alors que $v(N)>0$. Que la valeur de Shapley soit la seule règle à avoir ces propriétés et
l'additivité entre jeux, c'est le théorème de Shapley (1953), cité dans \texttt{omsources}.
\end{proof}

\omcode{code/firm/projsel/firm_projsel.py}{84}{108}{Contribution par exclusion, attribution selon l'ordre et valeur de Shapley sur un jeu $v$.}\label{lst:fm:managing-researchers-traders-and-engineers:shapley}

Dans le desk du chapitre (à titre d'illustration~: signaux A et B corrélés à 0.8, C indépendant, dix ans de données quotidiennes, le P\&L combiné ramené à 30
millions de dollars), les règles divergent complètement (la \cref{fig:fm:managing-researchers-traders-and-engineers:credit})~:

\begin{center}
\small
\begin{tabular}{@{}lrrrr@{}}
\toprule
contributeur & seul & \omterm{def:fm:managing-researchers-traders-and-engineers:credit}{contribution par exclusion} & ordre A, B, C & Shapley\\
\midrule
A & 22.77 & 2.46 & 22.77 & 12.63\\
B & 22.84 & 2.67 & 2.59 & 12.77\\
C & 4.55 & 4.64 & 4.64 & 4.61\\
somme & -- & 9.77 & 30.00 & 30.00\\
\bottomrule
\end{tabular}
\end{center}

\begin{omfigure}
\begin{tikzpicture}
\begin{axis}[width=10.5cm,height=5.2cm,ybar,bar width=8pt,ymin=0,ymax=25,ylabel={\small crédit (millions de \$)},
  xtick={0,1,2},xticklabels={A,B,C},tick label style={font=\footnotesize},table/col sep=comma,enlarge x limits=0.25,
  grid=major,grid style={gray!20},legend cell align=left,
  legend style={font=\scriptsize,at={(0.5,-0.25)},anchor=north,/tikz/every even column/.append style={column sep=8pt}},legend columns=4]
\addplot[fill=gray!40,draw=gray] table[x=k,y=alone]{figdata/desk/09-managing-researchers-traders-and-engineers/credit.csv};
\addplot[fill=omThm!50,draw=omThm] table[x=k,y=loo]{figdata/desk/09-managing-researchers-traders-and-engineers/credit.csv};
\addplot[fill=omProp!55,draw=omProp] table[x=k,y=order_abc]{figdata/desk/09-managing-researchers-traders-and-engineers/credit.csv};
\addplot[fill=omDef!55,draw=omDef] table[x=k,y=shapley]{figdata/desk/09-managing-researchers-traders-and-engineers/credit.csv};
\legend{seul,contribution par exclusion,{ordre A, B, C},Shapley}
\end{axis}
\end{tikzpicture}

\omcaption{Crédit d'une prévision combinée de 30 millions de dollars entre trois contributeurs selon quatre règles~: chaque signal seul, \omterm{def:fm:managing-researchers-traders-and-engineers:credit}{contribution par exclusion},
ordre d'arrivée (A en premier) et valeur de Shapley. A et B se recouvrent (corrélation 0.8), C est indépendant. Données illustratives. Source~:
\texttt{fm\_research.credit}.}\label{fig:fm:managing-researchers-traders-and-engineers:credit}
\end{omfigure}

Une enveloppe de bonus de 6~millions de dollars, répartie au prorata des attributions de chaque règle, donne à A, B et C 1.51, 1.64 et 2.85 millions de dollars selon la \omterm{def:fm:managing-researchers-traders-and-engineers:credit}{contribution par exclusion}, 2.53,
2.55 et 0.92 millions de dollars selon Shapley, et 4.55, 0.52 et 0.93 millions de dollars selon l'ordre d'arrivée. La \omterm{def:fm:managing-researchers-traders-and-engineers:credit}{contribution par exclusion} rémunère le plus C, car C est le seul
dont le signal ne peut être remplacé~; l'ordre d'arrivée rémunère celui qui est arrivé en premier. Shapley rémunère A et B conjointement pour l'effet qu'ils ont tous deux mis en évidence, et C
pour ce que seul C a mis en évidence. Les attributions sont elles-mêmes des estimations~: en refaire le calcul sur dix autres années de données (exercice 7) fait varier la répartition de plusieurs millions.

\section{Garder la recherche honnête}

Les règles d'\omterm{def:fm:managing-researchers-traders-and-engineers:credit}{attribution du mérite} façonnent la recherche autant qu'elles la mesurent. Rémunérés par \omterm{def:fm:managing-researchers-traders-and-engineers:credit}{contribution par exclusion}, les chercheurs évitent les travaux qui se recouvrent, y compris la réplication
qui exposerait un signal fallacieux~; rémunérés selon l'ordre d'arrivée, ils se précipitent pour arriver les premiers et thésaurisent~; rémunérés sur le P\&L autonome, ils redécouvrent ce que le
book détient déjà. Rémunéré selon Shapley, un chercheur qui trouve une seconde voie vers un effet que le desk exploite déjà est rémunéré pour une partie de celui-ci, ce qui
récompense la robustesse et la réplication. Trois autres règles gardent la recherche honnête quelle que soit la règle d'\omterm{def:fm:managing-researchers-traders-and-engineers:credit}{attribution du mérite}~: chaque essai figure dans le journal de recherche~;
les résultats sont présentés face à un test pré-enregistré~; et les personnes qui examinent un signal ne sont pas celles qui en sont rémunérées.

\section{Traders et ingénieurs~: des travaux différents, des retours différents}

Les chercheurs, les traders et les ingénieurs reçoivent un retour à des vitesses différentes, et les gérer de la même façon échoue pour les trois
(la \cref{fig:fm:managing-researchers-traders-and-engineers:loops}). Un trader voit le verdict du marché en quelques minutes et doit être jugé sur ses décisions,
non seulement sur le P\&L de la journée, dont le bruit le domine~: l'attribution du P\&L au risque pris et à l'exécution (livre~1, chapitre~7) est
l'outil. Le travail d'un chercheur ne porte ses fruits, si jamais, qu'après des mois de tests et des années de trading en conditions réelles~: la mesure est la qualité du processus, des essais
et des preuves, jusqu'à ce que le P\&L puisse parler. Le travail d'un ingénieur se voit dans des systèmes qui tournent ou non~: la livraison, la fiabilité et le temps nécessaire pour
rétablir le service se mesurent au quotidien (livre~15, chapitre~28), et la valeur du travail est le trading qu'il permet, ce qui explique pourquoi les ingénieurs
doivent figurer dans l'\omterm{def:fm:managing-researchers-traders-and-engineers:credit}{attribution du mérite} des stratégies qu'ils rendent possibles.

\begin{omfigure}
\begin{tikzpicture}[font=\footnotesize,>=stealth,box/.style={draw,rounded corners,minimum height=0.62cm,align=center,fill=omDef!10,
  minimum width=3.1cm}]
\node[box] (t) at (0,2.4) {trader\\{\scriptsize décision $\to$ marché~: minutes}};
\node[box] (e) at (0,1.2) {ingénieur\\{\scriptsize modification $\to$ production~: jours}};
\node[box] (r) at (0,0) {chercheur\\{\scriptsize idée $\to$ P\&L réel~: des mois à des années}};
\node[align=left,anchor=west,font=\scriptsize] at (2.8,2.4) {juger les décisions à l'aune de l'attribution~;\\le P\&L d'une journée est surtout du bruit};
\node[align=left,anchor=west,font=\scriptsize] at (2.8,1.2) {juger la livraison et la fiabilité~;\\attribuer le mérite du trading rendu possible};
\node[align=left,anchor=west,font=\scriptsize] at (2.8,0) {juger le processus et les preuves (essais, tests)\\jusqu'à ce que le P\&L puisse parler};
\end{tikzpicture}

\omcaption{Trois types de travail et la vitesse de leur retour d'information~: comment chacun doit être jugé tant que ses résultats ne sont encore que du bruit.}\label{fig:fm:managing-researchers-traders-and-engineers:loops}
\end{omfigure}

\section{Tutoriel~: une année de recherche et l'attribution du mérite}\label{tut:fm:managing-researchers-traders-and-engineers:tut}

\begin{tutorial}
\textbf{Objectif.} Choisir un \omterm{def:fm:managing-researchers-traders-and-engineers:portfolio}{portefeuille de recherche}, déflater un résultat par son nombre d'essais, et répartir le mérite d'une prévision combinée. \textbf{État final~:} le
tableau des attributions et la \cref{fig:fm:managing-researchers-traders-and-engineers:credit}.
\begin{tutsteps}
\item \textbf{Le portefeuille.} \texttt{firm.projsel.select(fm\_research.PROJECTS, 10)} choisit six projets~; \texttt{fm\_research.portfolio()}
simule 20\,000 années de ce choix et de choix aléatoires.
\item \textbf{Les essais.} \texttt{fm\_research.deflated(trials=40)} donne les ratios de Sharpe probabiliste et déflaté d'un backtest de cinq ans
avec un ratio de Sharpe de 1.2.
\item \textbf{Le jeu.} \texttt{firm.projsel.game(signals, target)} construit $v$~; \texttt{leave\_one\_out}, \texttt{ordered} et \texttt{shapley}
le répartissent (le \cref{lst:fm:managing-researchers-traders-and-engineers:shapley}).
\item \textbf{L'enveloppe de bonus.} \texttt{fm\_research.pool\_split()} convertit l'attribution de chaque règle en parts d'une enveloppe de bonus de 6~millions de dollars.
\end{tutsteps}
\end{tutorial}

\textbf{Pour aller plus loin.} Rendre A et B identiques, vérifier que Shapley les rémunère à égalité et que la \omterm{def:fm:managing-researchers-traders-and-engineers:credit}{contribution par exclusion} ne leur verse rien~; relancer avec
un autre germe aléatoire et mesurer l'incertitude de chaque attribution.

\section{Construction~: portefeuille de recherche et attribution du mérite}\label{bld:fm:managing-researchers-traders-and-engineers:projsel}

\begin{build}
\bfield{Rôle} Les deux problèmes d'allocation du responsable de desk~: quels projets financer, et comment attribuer le mérite aux personnes dont les travaux se sont combinés.
\bfield{Interface} \texttt{firm.projsel}~: \texttt{Project(name, cost, p, value, overlap)}, \texttt{expected\_net}~;
\texttt{select(projects, budget)}, \texttt{select\_random}, \texttt{simulate\_year}~; \texttt{game(signals, target)}~; \texttt{leave\_one\_out}, \texttt{ordered},
\texttt{shapley(v, n, samples, rng)}. Le chapitre utilise \texttt{firm.multitest} du livre~4 pour les ratios de Sharpe déflatés.
\bfield{Règles} Les projets à valeur nette espérée négative ou nulle ne sont jamais financés~; la valeur du jeu est le P\&L dans l'échantillon (in-sample) de la combinaison par
moindres carrés~; valeurs de Shapley exactes pour jusqu'à environ huit contributeurs, échantillonnées au-delà.
\bfield{Tests d'acceptation} \texttt{code/firm/projsel/tests/}~: sélection gloutonne dans le budget~; efficience de Shapley, symétrie entre
contributeurs identiques et nullité pour un contributeur factice~; \omterm{def:fm:managing-researchers-traders-and-engineers:credit}{contribution par exclusion} nulle pour un contributeur dupliqué~; valeur de Shapley échantillonnée proche de la valeur exacte.
\bfield{Pour aller plus loin} Jeux hors échantillon (out-of-sample), où la valeur de $S$ est mesurée sur des données non utilisées pour ajuster les pondérations~; un portefeuille dont les résultats de projets
sont corrélés~; \omterm{def:fm:managing-researchers-traders-and-engineers:credit}{attribution du mérite} avec une dimension temporelle (qui a maintenu un signal en vie).
\end{build}

\begin{omsources}
\item L.~S.~Shapley, «~A value for n-person games~», dans \emph{Contributions to the Theory of Games II}, 1953.
\item S.~Lipovetsky et M.~Conklin, «~Analysis of regression in game theory approach~», \emph{Applied Stochastic Models in Business and Industry}
17(4), 2001.
\item D.~H.~Bailey et M.~López de Prado, «~The deflated Sharpe ratio~», \emph{Journal of Portfolio Management} 40(5), 2014.
\item C.~R.~Harvey, Y.~Liu et H.~Zhu, «~\dots\ and the cross-section of expected returns~», \emph{Review of Financial Studies} 29(1), 2016.
\end{omsources}

\section{Exercices}

\begin{exercise}[$\star$]\label{exo:fm:managing-researchers-traders-and-engineers:1}
Le nouveau signal sur contrats à terme coûte 1.5~million de dollars, réussit avec une probabilité de 0.30 et vaut ensuite 12~millions de dollars. Calculer sa valeur nette espérée, ainsi que celle
rapportée à chaque dollar de coût~?
\end{exercise}

\begin{exercise}[$\star$]\label{exo:fm:managing-researchers-traders-and-engineers:2}
La deuxième variante de momentum coûte 1.0~million de dollars, réussit avec une probabilité de 0.40, vaut 6~millions de dollars et recouvre le livre à hauteur de 70~\%. Faut-il
la financer~?
\end{exercise}

\begin{exercise}[$\star$]\label{exo:fm:managing-researchers-traders-and-engineers:3}
Vérifier que les attributions de Shapley du tableau totalisent le P\&L combiné, puis calculer le montant que la \omterm{def:fm:managing-researchers-traders-and-engineers:credit}{contribution par exclusion} (leave-one-out) laisse non attribué.
\end{exercise}

\begin{exercise}[$\star\star$]\label{exo:fm:managing-researchers-traders-and-engineers:4}
Deux contributeurs apportent des signaux identiques valant 20 seuls et 20 ensemble~; un troisième apporte un signal indépendant valant 10. Calculer les attributions par
exclusion (leave-one-out) et de Shapley.
\end{exercise}

\begin{exercise}[$\star\star$]\label{exo:fm:managing-researchers-traders-and-engineers:5}
D'après la \cref{fig:fm:managing-researchers-traders-and-engineers:dsr}, au bout de combien de variantes le backtest du chapitre passe-t-il sous une chance sur deux
de battre le hasard~?
\end{exercise}

\begin{exercise}[$\star\star$]\label{exo:fm:managing-researchers-traders-and-engineers:6}
Quelle règle d'attribution récompense un chercheur qui reproduit un signal que le desk trade déjà, et pourquoi cela a-t-il de la valeur~?
\end{exercise}

\begin{exercise}[$\star\star\star$]\label{exo:fm:managing-researchers-traders-and-engineers:7}
\emph{Programmation.} Relancer \texttt{fm\_research.credit(seed=12)}. Comment les attributions par exclusion (leave-one-out) et de Shapley évoluent-elles, et que dit cela du
paiement sur l'attribution d'une seule année~?
\end{exercise}

\begin{exercise}[$\star\star\star$]\label{exo:fm:managing-researchers-traders-and-engineers:8}
\emph{Trouver la faille.} «~Nous versons à chaque chercheur 10~\% de ce que le desk perdrait si son signal était désactivé. C'est équitable~: c'est exactement
ce qu'il apporte.~»
\end{exercise}

\section{Problème~: qui a gagné l'argent}

\begin{problem}[{Problème du week-end --- qui a gagné l'argent}]\label{pb:fm:managing-researchers-traders-and-engineers:1}
Une prévision combinée a rapporté 30~millions de dollars. Trois chercheurs demandent leur part d'une enveloppe de 6~millions de dollars, et le responsable du desk doit aussi financer la recherche de l'année
suivante.

\medskip\noindent\textbf{Partie I --- Le portefeuille.}
\begin{enumerate}
\item Définir un \omterm{def:fm:managing-researchers-traders-and-engineers:portfolio}{portefeuille de recherche} et un \omterm{def:fm:managing-researchers-traders-and-engineers:portfolio}{tableau de bord de projet}.
\item Quels six projets la règle gloutonne finance-t-elle, et pour quelle valeur nette espérée~?
\item Donner la valeur réalisée moyenne et la probabilité d'une année perdante pour les choix gloutons et pour les choix aléatoires.
\item D'où viennent les probabilités de succès d'un \omterm{def:fm:managing-researchers-traders-and-engineers:portfolio}{tableau de bord de projet}, et pourquoi sont-elles optimistes~?
\end{enumerate}

\medskip\noindent\textbf{Partie II --- La revue.}
\begin{enumerate}[resume]
\item Qu'est-ce qui fait d'une revue de recherche une mise en scène, et qu'est-ce qui l'en empêche~?
\item Donner le ratio de Sharpe probabiliste et le ratio de Sharpe déflaté du backtest sur cinq ans après 40 variantes.
\item Quel est le meilleur ratio de Sharpe espéré de 40 variantes sans valeur~?
\item Pourquoi une nouvelle hypothèse doit-elle ouvrir un nouveau nombre d'essais~?
\end{enumerate}

\medskip\noindent\textbf{Partie III --- L'attribution.}
\begin{enumerate}[resume]
\item Définir la \omterm{def:fm:managing-researchers-traders-and-engineers:credit}{contribution par exclusion} (leave-one-out) et l'\omterm{def:fm:managing-researchers-traders-and-engineers:credit}{attribution du mérite}.
\item Donner l'attribution de chaque chercheur seul, par exclusion, selon l'ordre d'arrivée et selon Shapley.
\item Énoncer la \cref{prop:fm:managing-researchers-traders-and-engineers:shapley} et démontrer sa première partie.
\item Pourquoi la \omterm{def:fm:managing-researchers-traders-and-engineers:credit}{contribution par exclusion} laisse-t-elle ici l'essentiel du profit non attribué~?
\item Répartir l'enveloppe de 6~millions de dollars par exclusion, selon Shapley et selon l'ordre d'arrivée.
\item Quelle est l'incertitude de ces attributions~?
\end{enumerate}

\medskip\noindent\textbf{Partie IV --- Les personnes.}
\begin{enumerate}[resume]
\item Que pousse à faire chaque règle d'attribution chez les chercheurs~?
\item Comment juger les traders et les ingénieurs tant que leurs résultats ne sont que du bruit~?
\item Pourquoi les ingénieurs ont-ils leur place dans l'\omterm{def:fm:managing-researchers-traders-and-engineers:credit}{attribution du mérite} d'une stratégie~?
\item Qui doit revoir un signal, et qui ne le doit pas~?
\item Énoncer le \emph{résultat nommé}~: les attributions de Shapley et par exclusion des trois chercheurs, ainsi que la répartition de l'enveloppe qui en découle.
\item En deux phrases, expliquer aux trois chercheurs comment l'enveloppe a été répartie.
\end{enumerate}
\end{problem}

\section{Questions d'entretien}

\begin{interviewq}[$\star$ \iqroles{chercheur}]\label{iq:fm:managing-researchers-traders-and-engineers:1}
Deux de vos signaux sont fortement corrélés. Que devient la \omterm{def:fm:managing-researchers-traders-and-engineers:credit}{contribution par exclusion} (leave-one-out) de chacun à la prévision combinée~?
\end{interviewq}

\begin{interviewq}[$\star$ \iqroles{chercheur}]\label{iq:fm:managing-researchers-traders-and-engineers:2}
Votre backtest affiche un ratio de Sharpe de 1.2 sur cinq ans. Que devez-vous ajouter avant que quiconque puisse y croire~?
\end{interviewq}

\begin{interviewq}[$\star\star$ \iqroles{chercheur, trader}]\label{iq:fm:managing-researchers-traders-and-engineers:3}
Définissez la valeur de Shapley et donnez une propriété que l'attribution par exclusion ne possède pas.
\end{interviewq}

\begin{interviewq}[$\star\star$ \iqroles{développeur}]\label{iq:fm:managing-researchers-traders-and-engineers:4}
Comment mesureriez-vous la contribution d'une équipe d'ingénierie à un desk de trading~?
\end{interviewq}

\begin{interviewq}[$\star\star$ \iqroles{chercheur}]\label{iq:fm:managing-researchers-traders-and-engineers:5}
Vous disposez d'un budget pour cinq projets sur douze. Comment choisissez-vous, et quelle donnée d'entrée jugez-vous la moins fiable~?
\end{interviewq}

\begin{interviewq}[$\star\star\star$ \iqroles{chercheur}]\label{iq:fm:managing-researchers-traders-and-engineers:6}
Le calcul exact des valeurs de Shapley nécessite $n!$ ordres. Comment attribueriez-vous le mérite entre 30 contributeurs~?
\end{interviewq}
